\section{Conclusion}

This research proposal outlines a structured and empirically grounded plan to investigate how a learner-facing dashboard integrated into JupyterLab can support students' reflective practices and improve their programming processes. By combining fine-grained notebook interaction logs with both quantitative and qualitative measures of learning behavior, the study aims to provide novel evidence on how students engage with computational notebook environments beyond their final code outputs. The controlled experiment comparing students who use the dashboard with those who do not seeks to determine whether visualizing one's coding patterns and errors can foster metacognitive awareness, reduce repeated mistakes, and enhance learning strategies.

The expected outcomes of this work will contribute to the growing body of research on learning analytics in programming education by shifting the focus from course-level metrics to in-situ, process-oriented indicators. Additionally, the findings have the potential to inform the design of future dashboard systems and pedagogical interventions that make students' invisible learning processes more transparent to both learners and instructors. Ultimately, this project aims to advance our understanding of how reflective tools integrated directly within computational notebooks can support deeper learning and more effective problem-solving in introductory programming contexts.

% Next steps

Building on the proposed methodology, the next steps involve developing and carrying out the first initial survey to collect information about the participants' prior experiences and both their and their instructors' needs regarding the dashboard features. This information will then guide the design of the dashboard's metrics and visualizations to ensure they are relevant and actionable for the target users. Following this, and after the implementation of the dashboard extension and the recruitment of participants, the two learning units and associated data collection activities can be conducted.